\documentclass[11pt]{amsart}
\usepackage[T1]{fontenc}
\usepackage{lmodern}
\usepackage{amsmath,amssymb,amsthm,booktabs}
\usepackage{xcolor}
\usepackage[colorlinks=true,linkcolor=blue!55!black,citecolor=blue!55!black,urlcolor=blue!55!black]{hyperref}
\raggedbottom
\newtheorem{theorem}{Theorem}
\newtheorem{lemma}[theorem]{Lemma}
\newtheorem{proposition}[theorem]{Proposition}
\title[The finite covering invariant at 41]{The Exact Value of the Finite Erd\H{o}s--Ruzsa Covering Invariant at 41}
\author{Zijian Zeng}
\address{Institute of Computer Science and Digital Innovation, UCSI University, Kuala Lumpur, 56000, MALAYSIA}
\email{zijianzeng@foxmail.com}
\email{1002266693@ucsiuniversity.edu.my}
\date{September 1, 2026}
\hypersetup{pdftitle={The Exact Value of the Finite Erdos--Ruzsa Covering Invariant at 41},pdfauthor={Zijian Zeng}}
\begin{document}
\begin{abstract}
For a positive integer $N$, let $\mu(N)$ be the minimum of
$\sum_{p\in P}1/p$ over subsets $P$ of the primes at most $N$ for which one
residue class modulo each $p\in P$ covers $[1,N]$.  We prove the certified
finite result
\[
 \mu(41)=\frac{4369938619315}{3202634352918}.
\]
An explicit cover proves the upper bound.  The lower bound is a complete exact
enumeration using forced branching through an uncovered integer and an exact
fractional-capacity relaxation.  No floating-point calculation or optimizer
status enters the proof.  This finite result does not resolve the asymptotic
Erd\H{o}s--Ruzsa covering problem.
\end{abstract}
\maketitle
\section{Introduction}
Erd\H{o}s and Ruzsa conjectured that, for all sufficiently large $x$, the positive integers
below $x$ can be covered by one residue class for each prime in a selected
finite set, while the sum of the reciprocal moduli remains bounded by an
absolute constant~\cite[p.~106]{Erdos1980}.  Erd\H{o}s and Ruzsa also posed the
complementary arbitrary-residue small-sieve question~\cite[p.~386, Problem~2]{ErdosRuzsa1980}.
The asymptotic problem is currently recorded as open~\cite{Bloom1200}.

The direct finite invariant below converts the question into exact weighted
set cover.  A public computational working report dated July 27, 2026,
reported a complete table through $N=40$, ending with~\cite{WorkingReport1200}
\[
\mu(40)=\frac{2067007816901}{1484147626962}.
\]
We determine the next cell.  Our novelty search covered the cited primary
sources, the live problem record, exact-value web searches, and OEIS through
September 1, 2026; no earlier value for $\mu(41)$ was located.  This is a
date-bounded search report, not a claim that an unindexed computation cannot
exist.

\section{Definition and exact normalization}
For $N\ge1$, write $\mathcal P_N=\{p\le N:p\text{ is prime}\}$ and define
\[
\mu(N)=\min\left\{\sum_{p\in P}\frac1p:P\subseteq\mathcal P_N
\text{ admits a covering residue assignment}\right\}.
\]
Here a covering residue assignment means integers $0\le a_p<p$, one for each
$p\in P$, such that for every $n\in[1,N]$ there is a $p\in P$ with
$n\equiv a_p\pmod p$.
The value is $+\infty$ if no such cover exists.  Thus primes may be omitted,
but at most one class is used for each selected prime.  For $N=41$, every
objective comparison may be multiplied by
\[
 L=\prod_{p\le41}p=304250263527210,
\]
turning it into an integer comparison.

\section{The exact value}
\begin{theorem}\label{thm:main}
One has
\[
\boxed{\displaystyle \mu(41)=\frac{4369938619315}{3202634352918}}.
\]
\end{theorem}
\begin{proof}
For the upper bound, use the following classes.
\[
\begin{array}{c|rrrrrrrrrrr}
p&2&3&7&11&13&17&23&29&31&37&41\\ \midrule
a_p&1&1&3&8&6&2&18&14&20&12&26
\end{array}
\]
Direct reduction modulo the displayed primes shows that every integer from 1
through 41 belongs to at least one class.  Summing the reciprocals and reducing
gives
\[
 \frac12+\frac13+\frac17+\frac1{11}+\frac1{13}+\frac1{17}
 +\frac1{23}+\frac1{29}+\frac1{31}+\frac1{37}+\frac1{41}
 =\frac{4369938619315}{3202634352918}.
\]
Proposition~\ref{prop:enumeration} proves that no cheaper cover exists.
\end{proof}

\section{Complete lower-bound certificate}
Represent a subset of $[1,41]$ by a 41-bit integer.  For each prime $p\le41$
and $0\le a<p$, precompute the mask
$M_{p,a}=\{n\in[1,41]:n\equiv a\pmod p\}$.  A state consists of a covered mask
$C$ and the set $R$ of unused primes.
\begin{lemma}[Complete forced branching]\label{lem:branch}
If $n\notin C$, every completion of the state uses, for some $p\in R$, the
unique class $M_{p,n\bmod p}$.
\end{lemma}
\begin{proof}
Any completion covers $n$ with a class modulo an unused prime $p$.  Exactly one
residue modulo $p$ contains $n$, namely $n\bmod p$.
\end{proof}

For pruning, let $U=[1,41]\setminus C$ and
\[
 g_p=\max_{0\le a<p}|U\cap M_{p,a}|.
\]
Relax every remaining prime to a divisible item of capacity $g_p$ and cost
$1/p$, allow fractional items, and ignore overlaps.  Filling $|U|$ units in
nondecreasing order of $(1/p)/g_p$ gives a lower bound on every completion:
the relaxation has weakly more freedom than the residue-class problem.  The
implementation evaluates this bound after multiplication by $L$, using integer
cross multiplication even for the final fractional item.
\begin{proposition}\label{prop:enumeration}
There is no cover of $[1,41]$ having cost strictly less than
$4369938619315/3202634352918$.
\end{proposition}
\begin{proof}
Starting from $C=\varnothing$ and all 13 primes unused, choose an uncovered
$n$ and branch over every possibility in Lemma~\ref{lem:branch}.  Discard a
state only if its accumulated exact integer cost already reaches the strict
budget, or if the fractional-capacity lower bound reaches it.  Cache only
states whose entire descendant tree has failed.  The pair $(C,R)$ determines
the selected prime set and hence its accumulated cost, so this caching loses no
budget information.

The completed run entered 28,998,666 states, cached 5,536,355 failed states,
and applied 22,110,683 capacity prunes.  It returned no cheaper cover.  The
source recomputes the primes, masks, exact objective, and incumbent coverage;
the serialized run record includes the interpreter and platform versions.
By Lemma~\ref{lem:branch}, every possible completion occurs on a branch, and
the only nontrivial prune is a lower bound from a relaxation.  The enumeration
is therefore exhaustive.
\end{proof}

As a structurally independent small-instance check, a Cartesian program
enumerated, for every prime, the choice ``unused'' or every residue.  For all
$1\le N\le10$, its finite optima and infeasible cells agree with the main
enumerator; in particular the finite values are $5/6$ at $N=3$, $31/30$ at
$N=5$, and $247/210$ at $N=7,8,9$.

\section{Scope and reproducibility}
The theorem is a certified finite result only.  In particular, it neither
proves nor disproves boundedness of $\mu(N)$ as $N\to\infty$.  The exact
witness, exhaustive enumerator, independent witness checker, small-instance
Cartesian checker, logs, and commands accompany the source archive.  The fast
checker recomputes primality, residue validity, coverage, and the reduced
rational cost from the serialized witness.  The full certificate uses Python
3.10.12 and only its standard library; it is deterministic and uses no random
seed.  Optimization software was used to locate a candidate, but its status
and floating-point objective are not used in the proof.
\bibliographystyle{amsplain}
\bibliography{references}
\end{document}
